% Anexo: matemática del método VBF (Voigt-Based Fitting, Rancourt–Ping 1991).
% Se incluye con \input desde main.tex bajo un \chapter, así que empieza en \section.

\section{Motivation: a basis of Gaussians for \texorpdfstring{$P(\BHF)$}{P(BHF)}}
\label{sec:vbf-motiv}

In a magnetically disordered solid —a metallic glass, a substituted ferrite, a
nanoparticle, an amorphous material— the hyperfine field $\BHF$ seen by each
$^{57}$Fe nucleus is not unique, but is distributed according to a probability
density $P(\BHF)$. The measured spectrum is the \emph{superposition} of all the
elemental sextets weighted by that distribution. Recovering $P(\BHF)$ from the
spectrum is an ill-conditioned inverse problem (see the appendix of the
mathematical manual, \S\ref{sec:dist-corr} and nearby sections).

The Hesse--Rübartsch \textbf{histogram} solves it without assuming a shape: it
assigns a free weight to each bin and controls the oscillations with a
regularizer (Tikhonov, TV or MaxEnt) governed by $\alpha$. It is flexible, but it
delivers a profile without structure and is sensitive to the choice of $\alpha$.

The \textbf{VBF} (\emph{Voigt-Based Fitting}, Rancourt and Ping, 1991) adopts the
opposite hypothesis, which is very reasonable in practice: each magnetic
\textbf{environment} or \textbf{subpopulation} contributes a \emph{Gaussian} of
fields, and the total distribution is the sum of a few of them,
\begin{equation}
  P(\BHF) \;=\; \sum_{j=1}^{N} A_j\,\mathcal N\!\bigl(\BHF;\,\mu_j,\sigma_j\bigr),
  \qquad N = 1\ldots6 .
  \label{eq:vbf-P}
\end{equation}
Instead of dozens of bin weights, the model has $3N$ shape parameters with
\textbf{direct physical meaning}: the position $\mu_j$ (mean field of the
environment), the width $\sigma_j$ (its dispersion) and the area $A_j$ (its
population). By reducing the degrees of freedom so much, the fit is stable and
\textbf{does not require} regularization or $\alpha$.

\section{From field to spectrum: the forward model}
\label{sec:vbf-directo}

Fitbauer discretizes the variable over a grid of $M$ nodes
$c_1<\dots<c_M$ (the bin centers, \code{parameter\_grid}). To each node $c_k$
corresponds an \textbf{elemental subspectrum} —an Fe-57 sextet with field
$c_k$, shift $\IS$, quadrupole splitting $\dEQ$ and width $\gamma$—
whose absorption in the velocity channel $v_i$ is the element of the
\textbf{kernel}
\begin{equation}
  K_{ik} \;=\; \sum_{\ell=1}^{6} I_\ell(c_k)\;
  \Phi\!\bigl(v_i - p_\ell(c_k);\,\Gamma_\ell,\varsigma\bigr),
  \label{eq:vbf-kernel}
\end{equation}
where $p_\ell(c_k)$ and $I_\ell(c_k)$ are the positions and intensities of the six
lines (appendix~\ref{ap:mates}) and $\Phi$ is the line profile
(\S\ref{sec:vbf-voigt}). Over this grid, equation~\eqref{eq:vbf-P} becomes the
weight vector
\begin{equation}
  w_k \;=\; \sum_{j=1}^{N} A_j\, \mathcal N(c_k;\mu_j,\sigma_j),
  \qquad k=1,\dots,M,
  \label{eq:vbf-w}
\end{equation}
and the modeled spectrum is linear in those weights:
\begin{equation}
  \hat y_i \;=\; b + s\,v_i \;-\; \sum_{k=1}^{M} K_{ik}\,w_k
  \;-\; \underbrace{\sum_{m} a_m\,S_{im}}_{\text{sharp components (opt.)}},
  \label{eq:vbf-model}
\end{equation}
with baseline $b$, slope $s$ and, optionally, one or several \emph{sharp}
subspectra $S_{im}$ (discrete crystalline phases) of amplitude $a_m\ge0$.

\paragraph{Gaussian kernel normalized to unit area.} The elemental Gaussian of
\eqref{eq:vbf-w} is evaluated and \textbf{normalized to unit area} over the grid
itself,
\begin{equation}
  \mathcal N(c_k;\mu,\sigma) \;=\;
  \frac{\exp\!\bigl(-\tfrac12((c_k-\mu)/\sigma)^2\bigr)}
       {\displaystyle\int \exp\!\bigl(-\tfrac12((c-\mu)/\sigma)^2\bigr)\,\dd c},
  \label{eq:vbf-gauss}
\end{equation}
where the denominator is computed by the trapezoidal rule (\code{gauss\_weights}).
This normalization is what gives the coefficient $A_j$ the meaning of \emph{area}
(\S\ref{sec:vbf-areas2}).

\section{Why "Voigt-based"}
\label{sec:vbf-voigt}

The name of the method comes from the line profile $\Phi$. The intrinsic
Mössbauer line is \textbf{Lorentzian} of natural width, but the instrumental
broadening and the finite thickness of the absorber convolve it with a
\textbf{Gaussian}. The convolution of the two is the \textbf{Voigt} profile
\begin{equation}
  V(x;\gamma,\varsigma) \;=\; (L_\gamma * G_\varsigma)(x)
  \;=\; \frac{\operatorname{Re}\,\mathrm w(z)}{\varsigma\sqrt{2\pi}},
  \qquad
  z = \frac{x + i\,\gamma/2}{\varsigma\sqrt 2},
  \label{eq:vbf-voigt}
\end{equation}
with $L_\gamma$ a Lorentzian of full width $\gamma$, $G_\varsigma$ a Gaussian of
deviation $\varsigma$ (the \param{Voigt $\sigma$} parameter of the GUI) and
$\mathrm w(z)$ the Faddeeva function. As $\varsigma\to0$ the pure Lorentzian is
recovered. In \eqref{eq:vbf-kernel} one takes $\Phi=V$ ("Voigt" profile) or
$\Phi=L$ ("Lorentz" profile); Fitbauer sets the profile explicitly when building
the VBF kernel.

\paragraph{The dual reading of the Gaussian broadening.} There is an important
equivalence that justifies the approach. Distributing a \emph{Lorentzian} sextet
over a \emph{Gaussian of fields} of width $\sigma_j$ is equivalent, on the
velocity axis, to convolving that sextet with a Gaussian whose width is
$\sigma_j$ scaled by the sensitivity $\partial p_\ell/\partial\BHF$ of each line
to the field. That is, \textbf{each component of $P(\BHF)$ produces a sextet of
effectively \emph{Voigtian} lines}: hence "Voigt-Based Fitting". The parameter
$\sigma_j$ of the distribution and the broadening $\varsigma$ of the profile
capture, respectively, the \emph{physical} field dispersion of the environment and
the \emph{instrumental} broadening; Fitbauer keeps them separate so as not to
confuse one with the other.

\begin{figure}[htbp]\centering
  \includegraphics[width=0.82\textwidth]{apx_vbf}
  \caption{VBF decomposition (Rancourt--Ping). The distribution $P(\BHF)$ is
  modeled as a sum of Gaussians \eqref{eq:vbf-P}; since the Gaussian kernel is
  normalized to unit area \eqref{eq:vbf-gauss}, the amplitude $A_j$ of each
  component \emph{is} its area and its fraction $A_j/\sum_i A_i$ gives the
  population of the environment. The sum reproduces the distribution with few
  parameters and without regularization.}
  \label{fig:vbf-decomp}
\end{figure}

\section{Least-squares problem and parametrization}
\label{sec:vbf-lsq}

The fit minimizes the residual of \eqref{eq:vbf-model} by nonlinear least squares
(\code{scipy.optimize.least\_squares}, trust region) over the parameter vector
\begin{equation}
  \bm x \;=\; \bigl(\,b,\; s,\; \{A_j,\ \mu_j,\ \log\sigma_j\}_{j=1}^{N},\;
  \{a_m\}_m\,\bigr).
  \label{eq:vbf-x}
\end{equation}
Details of the parametrization:
\begin{itemize}[nosep]
  \item \textbf{$\log\sigma_j$ is fitted}, not $\sigma_j$: it guarantees
        $\sigma_j>0$ without a hard barrier and balances the scale of the width
        against the position and the area, improving the conditioning.
  \item \textbf{Physical bounds}: $A_j\ge0$ (non-negative areas),
        $\mu_j\in[\,B_{\mathrm{hf},\min},B_{\mathrm{hf},\max}]$ within the grid, and
        $\sigma_j$ bounded between a fraction of a bin and the full range. The
        baseline satisfies $b\ge0$ and the sharp amplitudes $a_m\ge0$.
  \item \textbf{Start}: the $N$ centers are seeded equispaced over the range,
        with initial width $\sim(B_{\mathrm{hf},\max}-B_{\mathrm{hf},\min})/(3N)$ and area
        distributed.
  \item \textbf{Ordering and identifiability}: since the model is invariant under
        permutations of the components (\emph{label switching}), on completion they
        are \textbf{ordered by $\mu_j$} before being saved, so that component $1$
        is always the one with the lowest field.
\end{itemize}

Compared with the histogram, here there is \emph{no} regularization term: the
functional form \eqref{eq:vbf-P} itself imposes the smoothness, and for this reason
the selection of $\alpha$ and the L-curve do not intervene (the corresponding
button is disabled).

\section{Areas, populations and percentages}
\label{sec:vbf-areas2}

Since each Gaussian kernel of \eqref{eq:vbf-gauss} integrates to unity over the
grid, the area of component $j$ in $P(\BHF)$ is exactly its coefficient:
\begin{equation}
  \int A_j\,\mathcal N(\BHF;\mu_j,\sigma_j)\,\dd\BHF \;=\; A_j .
\end{equation}
Therefore the \textbf{population fraction} of environment $j$ —what proportion of
the Fe nuclei see that mean field— is, in percent,
\begin{equation}
  f_j \;=\; 100\;\frac{A_j}{\sum_{i=1}^{N} A_i}.
  \label{eq:vbf-pop}
\end{equation}
Fitbauer reports $A_j$, $f_j$, $\mu_j$ and $\sigma_j$ per component in the status
panel and in the reports. This is the operational advantage of the VBF: it
\textbf{directly quantifies the proportion} of two or more magnetic
subpopulations, something that a histogram only gives after integrating regions of
the profile by hand.

\section{Correlation \texorpdfstring{$\IS(H)$}{d(H)} and \texorpdfstring{$\dEQ(H)$}{DEQ(H)}}
\label{sec:vbf-corr}

The VBF admits the same correlation slopes as the histogram
(\S\ref{sec:dist-corr}): when building each column $k$ of the kernel
\eqref{eq:vbf-kernel} one uses
\begin{equation}
  \IS_k = \IS_0 + \kappa_\delta\,(c_k - H_{\text{ref}}),\qquad
  \Delta E_{Q,k} = \Delta E_{Q,0} + \kappa_q\,(c_k - H_{\text{ref}}),
\end{equation}
instead of global values (Le~Caër--Dubois / Wivel--M\o{}rup model). The model
remains linear in the weights $w_k$; only the kernel changes. With
$\kappa_\delta=\kappa_q=0$ (by default) the classic case is recovered, and the
slopes can be fixed or refined in the outer loop.

\section{Case \texorpdfstring{$N=1$}{N=1}: the "Gaussian" shape}
\label{sec:vbf-gaussiana}

The \menu{Gaussian} shape of the interface \textbf{is the VBF with a single
component} ($N=1$) and a Lorentzian line profile. With this, \eqref{eq:vbf-P}
reduces to a single Gaussian of fields,
\begin{equation}
  P(\BHF) = A\,\mathcal N(\BHF;\mu,\sigma),
\end{equation}
and the parameter vector \eqref{eq:vbf-x} to $(b,s,A,\mu,\log\sigma)$: the minimal
description of \emph{one} broadened, approximately symmetric environment.
Internally Fitbauer does not maintain a separate algorithm for it —"Gaussian" and
"VBF" share exactly the same code— but routes the label to the same fit with
$N=1$. Choose "Gaussian" as a quick first approximation to a single phase; move up
to "VBF" with $N=2$--$3$ as soon as the residual reveals more than one
environment.

\section{VBF versus the histogram}
\label{sec:vbf-vs-hist}

\begin{center}
\small
\begin{tabular}{@{}p{0.30\textwidth}p{0.32\textwidth}p{0.30\textwidth}@{}}
\toprule
 & \textbf{Histogram (H.--R.)} & \textbf{VBF} \\
\midrule
Shape parameters & $M$ bin weights ($\sim$40--100) & $3N$ ($N=1$--$6$) \\
Smoothness & imposed by $\alpha$ (regularizer) & imposed by the shape \\
Selection of $\alpha$ / L-curve & yes, critical & not applicable \\
Output & free profile & $A_j,\mu_j,\sigma_j$ per environment \\
Populations & integrating by hand & direct \eqref{eq:vbf-pop} \\
Risk & oscillations if $\alpha$ is wrong & overfitting if $N$ is high \\
\bottomrule
\end{tabular}
\end{center}

In summary: use the \textbf{histogram} when you do not want to commit to a shape
(a possibly complex or irregularly multimodal distribution) and the \textbf{VBF}
when you expect \emph{few well-defined environments} and are interested in their
position, width and proportion. Both share the same kernel, the same correlations
and the same sharp components.

\section{Best practices}
\label{sec:vbf-practica}

\begin{itemize}[nosep]
  \item \textbf{Start low}: $N=1$--$2$ and increase only if the residual demands
        it. With an excessive $N$, two Gaussians tend to overlap on the same
        environment and end up poorly determined (anticorrelated areas, degenerate
        widths).
  \item \textbf{Watch the $\kappa_\delta$--grid-shift degeneracy}: a slope
        $\kappa_\delta$ and a global shift of $\IS$ are almost indistinguishable;
        refine with bounds and gently.
  \item \textbf{Separate physics and instrument}: let $\sigma_j$ absorb the field
        dispersion and $\varsigma$ (Voigt) the instrumental broadening; do not
        compensate one with the other.
  \item \textbf{Compare with the histogram}: if the $P(\BHF)$ of the regularized
        histogram and the envelope of the VBF coincide, the description in a few
        Gaussians is justified.
\end{itemize}

\section*{References}
\begin{itemize}[nosep]
  \item D.~G.~Rancourt, J.~Y.~Ping (1991). \emph{Voigt-based methods for
        arbitrary-shape static hyperfine parameter distributions in Mössbauer
        spectroscopy}. Nucl. Instrum. Methods Phys. Res. B \textbf{58}, 85--97.
  \item G.~Le~Caër, J.~M.~Dubois (1979). \emph{Evaluation of hyperfine parameter
        distributions from overlapped Mössbauer spectra of amorphous alloys}.
        J. Phys. E: Sci. Instrum. \textbf{12}, 1083--1090.
  \item C.~Wivel, S.~M\o{}rup (1981). \emph{Improved computational procedure for
        evaluation of overlapping hyperfine parameter distributions in Mössbauer
        spectra}. J. Phys. E: Sci. Instrum. \textbf{14}, 605--610.
\end{itemize}
